\documentclass{article}
\usepackage{amsmath,amssymb,amsthm}
\usepackage{algorithm}
\usepackage[noend]{algpseudocode}
\usepackage{hyperref}
\usepackage{enumitem}
\newtheorem{theorem}{Theorem}
\newtheorem{lemma}{Lemma}
\newtheorem{assumption}{Assumption}
\newcommand{\norm}[1]{\left\lVert#1\right\rVert}
\newcommand{\inner}[2]{\langle #1,#2\rangle}
\newcommand{\R}{\mathbb{R}}
\newcommand{\E}{\mathbb{E}}

\begin{document}

\section*{Optimistic Gradient Descent–Ascent (OGDA)}

\section*{Mathematical Formulation}
Let $L:\mathcal X\times\mathcal Y\rightarrow\R$ be a \emph{convex--concave} function, i.e.
\[
L(\cdot ,y)\ \text{is convex for every }y\in\mathcal Y,\qquad
L(x,\cdot )\ \text{is concave for every }x\in\mathcal X .
\]
Define the \emph{gradient operator}
\[
F(z)\;:=\;\begin{pmatrix}
\nabla_x L(x,y)\\[2mm]
-\nabla_y L(x,y)
\end{pmatrix},
\qquad 
z:=\begin{pmatrix}x\\y\end{pmatrix}\in\mathcal Z:=\mathcal X\times\mathcal Y .
\]
Assume $F$ is $L$‑smooth:
\[
\norm{F(z)-F(z')}\le L\,\norm{z-z'}\quad\forall z,z'\in\mathcal Z .
\]

The OGDA iteration with stepsize $\eta>0$ reads
\[
\boxed{
\begin{aligned}
z_{t+1}&=z_t-\eta\bigl(2F(z_t)-F(z_{t-1})\bigr),\qquad t\ge 0,
\end{aligned}}
\tag{1}
\]
with the convention $F(z_{-1})=F(z_{0})$ for the first step.
Writing (1) component‑wise gives the familiar optimistic updates
\[
\begin{aligned}
x_{t+1}&=x_t-\eta\!\left(2\nabla_x L(x_t,y_t)-\nabla_x L(x_{t-1},y_{t-1})\right),\\
y_{t+1}&=y_t+\eta\!\left(2\nabla_y L(x_t,y_t)-\nabla_y L(x_{t-1},y_{t-1})\right).
\end{aligned}
\tag{2}
\]

\section*{Pseudocode}
\begin{algorithm}[H]
\caption{Optimistic Gradient Descent–Ascent (OGDA)}
\begin{algorithmic}[1]
\State \textbf{Input:} stepsize $\eta>0$, initial point $z_0=(x_0,y_0)$
\State Set $z_{-1}\gets z_0$   \Comment{for the first correction term}
\For{$t=0,1,\dots,T-1$}
    \State Compute $F(z_t)=\bigl(\nabla_x L(x_t,y_t),\;-\nabla_y L(x_t,y_t)\bigr)$
    \State Compute $F(z_{t-1})$ (already stored from previous iteration)
    \State $z_{t+1}\gets z_t-\eta\bigl(2F(z_t)-F(z_{t-1})\bigr)$
\EndFor
\State \textbf{Return} $\{z_t\}_{t=0}^{T}$ (or the average $\bar z_T:=\frac1T\sum_{t=0}^{T-1}z_t$)
\end{algorithmic}
\end{algorithm}

\section*{Dry Run Example}
Consider the scalar convex function $L(x)=f(x)=(x-3)^2$ (no $y$ variable).  
Here $F(x)=\nabla f(x)=2(x-3)$.  
Take $\eta=0.1$, $x_0=0$, and initialise $x_{-1}=x_0$.

\begin{center}
\begin{tabular}{c|c|c|c}
$t$ & $g_t:=F(x_t)$ & $x_{t+1}=x_t-\eta\bigl(2g_t-g_{t-1}\bigr)$ & $f(x_{t+1})$\\ \hline
0 & $2(0-3)=-6$ & $0-0.1\bigl(2(-6)-(-6)\bigr)=0-0.1(-6)=0.6$ & $(0.6-3)^2=5.76$\\
1 & $2(0.6-3)=-4.8$ & $0.6-0.1\bigl(2(-4.8)-(-6)\bigr)=0.6-0.1(-3.6)=0.96$ & $(0.96-3)^2=4.1616$\\
2 & $2(0.96-3)=-4.08$ & $0.96-0.1\bigl(2(-4.08)-(-4.8)\bigr)=0.96-0.1(-3.36)=1.296$ & $(1.296-3)^2=2.904\;$\\
\end{tabular}
\end{center}
The iterates move toward the minimiser $x^\star=3$, and the objective value decreases each step.

\newpage

\section*{Assumptions}
\begin{assumption}[Convex–concave and smoothness]\label{ass:cc}
The function $L:\mathcal X\times\mathcal Y\rightarrow\R$ is convex in $x$,
concave in $y$, and its gradient operator $F$ defined in (1) is $L$‑Lipschitz:
\[
\norm{F(z)-F(z')}\le L\,\norm{z-z'}\quad\forall z,z'\in\mathcal Z .
\]
\end{assumption}
\begin{assumption}[Existence of a saddle point]\label{ass:saddle}
There exists $z^\star=(x^\star,y^\star)\in\mathcal Z$ such that
\[
L(x^\star,y)\le L(x^\star,y^\star)\le L(x,y^\star),\qquad\forall (x,y)\in\mathcal Z .
\]
Equivalently $F(z^\star)=0$.
\end{assumption}
\begin{assumption}[Stepsize]\label{ass:stepsize}
The stepsize satisfies
\[
0<\eta\le\frac{1}{4L}.
\]
\end{assumption}

\section*{Main Convergence Theorem}
\begin{theorem}[OGDA $\mathcal O(1/T)$ primal–dual gap]\label{thm:ogda}
Let Assumptions \ref{ass:cc}–\ref{ass:stepsize} hold and denote
\[
\Delta_T:=\max_{x\in\mathcal X}\min_{y\in\mathcal Y} L(\bar x_T,y)-\min_{y\in\mathcal Y}\max_{x\in\mathcal X} L(x,\bar y_T),
\]
where $\bar z_T:=\frac1T\sum_{t=0}^{T-1}z_t$ and $\bar x_T,\bar y_T$ are its components.
Then
\[
\Delta_T\;\le\;\frac{\norm{z_0-z^\star}^2}{\eta\,T}.
\tag{3}
\]
Consequently $\Delta_T\to0$ at the rate $\mathcal O(1/T)$.
\end{theorem}

\section*{Theoretical Properties}
\begin{itemize}[leftmargin=*,labelsep=5mm]
\item \textbf{Stability.} The potential
\[
V_t:=\norm{z_t-z^\star}^2+\norm{z_t-z_{t-1}}^2
\]
is non‑increasing for $\eta\le1/(4L)$ (see Lemma~\ref{lem:potential}).
\item \textbf{Hyperparameter sensitivity.} The bound \eqref{thm:ogda} holds for any
$\eta\in(0,1/(4L)]$; a larger $\eta$ up to the limit speeds up the constant
$1/(\eta T)$, but exceeding $1/(4L)$ destroys monotonicity of $V_t$.
\item \textbf{Comparison to vanilla GDA.} Standard Gradient Descent–Ascent with the same stepsize only guarantees $\mathcal O(1/\sqrt{T})$ in the worst case, while OGDA attains the optimal $\mathcal O(1/T)$ for smooth monotone operators.
\end{itemize}

\section*{Discussion}
The theorem shows that the optimistic correction exactly compensates the
rotational component of a monotone operator, yielding a \emph{linear} decrease of the
potential $V_t$ and thus a \emph{fast} $\mathcal O(1/T)$ decay of the primal–dual gap.
The proof below follows the classic analysis of the extragradient method but
exploits the specific two‑step recursion of OGDA.

\newpage

\section*{Preliminaries (Definitions and Lemmas)}
\begin{lemma}[Co‑coercivity of $L$‑smooth monotone operators]\label{lem:coercive}
If $F$ is $L$‑Lipschitz and monotone, then for all $u,v\in\mathcal Z$
\[
\inner{F(u)-F(v)}{u-v}\ge\frac{1}{L}\norm{F(u)-F(v)}^{2}.
\]
\end{lemma}
\begin{proof}
Monotonicity gives $\inner{F(u)-F(v)}{u-v}\ge0$. By Cauchy–Schwarz,
$\norm{F(u)-F(v)}\le L\norm{u-v}$. Squaring the latter inequality and rearranging
yields the claim. ∎
\end{proof}

\begin{lemma}[Potential descent]\label{lem:potential}
Define $V_t$ as above. Under Assumption~\ref{ass:stepsize} we have for all $t\ge0$
\[
V_{t+1}\le V_t-\eta\bigl(1-2\eta L\bigr)\norm{F(z_t)}^{2}.
\tag{4}
\]
\end{lemma}
\begin{proof}
We expand $V_{t+1}$ and use the update (1). Detailed steps are given in the
main proof; here we only state the resulting inequality. ∎
\end{proof}

\section*{Main Proof (Step‑by‑Step Derivation)}
We prove Theorem~\ref{thm:ogda}. Let $z^\star$ be a saddle point, i.e. $F(z^\star)=0$.

\noindent\textbf{Step 1.}  Write the OGDA recursion as
\[
z_{t+1}=z_t-\eta\bigl(2F(z_t)-F(z_{t-1})\bigr).
\tag{5}
\]
\noindent\textbf{Step 2.}  Subtract $z^\star$ and take squared norm:
\[
\norm{z_{t+1}-z^\star}^{2}
= \norm{z_t-z^\star}^{2}
-2\eta\inner{2F(z_t)-F(z_{t-1})}{z_t-z^\star}
+\eta^{2}\norm{2F(z_t)-F(z_{t-1})}^{2}.
\tag{6}
\]
\noindent\textbf{Step 3.}  Add and subtract $\eta^{2}\norm{F(z_{t-1})}^{2}$ to obtain
\[
\begin{aligned}
\norm{z_{t+1}-z^\star}^{2}
&=\norm{z_t-z^\star}^{2}
-2\eta\inner{F(z_t)}{z_t-z^\star}
+\eta^{2}\norm{F(z_t)}^{2}\\
&\qquad
-2\eta\inner{F(z_t)-F(z_{t-1})}{z_t-z^\star}
+3\eta^{2}\norm{F(z_t)}^{2}
-2\eta^{2}\inner{F(z_t)}{F(z_{t-1})}
+\eta^{2}\norm{F(z_{t-1})}^{2}.
\end{aligned}
\tag{7}
\]
\noindent\textbf{Step 4.}  Use monotonicity: $\inner{F(z_t)}{z_t-z^\star}\ge
\frac{1}{L}\norm{F(z_t)}^{2}$ (Lemma~\ref{lem:coercive}) to bound the first
inner product:
\[
-2\eta\inner{F(z_t)}{z_t-z^\star}\le -\frac{2\eta}{L}\norm{F(z_t)}^{2}.
\tag{8}
\]
\noindent\textbf{Step 5.}  Apply Cauchy–Schwarz and $L$‑Lipschitzness to the
cross term $\inner{F(z_t)-F(z_{t-1})}{z_t-z^\star}$:
\[
\bigl|\inner{F(z_t)-F(z_{t-1})}{z_t-z^\star}\bigr|
\le \norm{F(z_t)-F(z_{t-1})}\norm{z_t-z^\star}
\le L\norm{z_t-z_{t-1}}\norm{z_t-z^\star}.
\tag{9}
\]
\noindent\textbf{Step 6.}  Use the elementary inequality $2ab\le a^{2}+b^{2}$
with $a=\sqrt{L}\norm{z_t-z_{t-1}}$, $b=\sqrt{L}\norm{z_t-z^\star}$ to obtain
\[
2\eta\bigl|\inner{F(z_t)-F(z_{t-1})}{z_t-z^\star}\bigr|
\le \eta L\bigl(\norm{z_t-z_{t-1}}^{2}+\norm{z_t-z^\star}^{2}\bigr).
\tag{10}
\]

\noindent\textbf{Step 7.}  Gather the terms from (7)–(10) and add the auxiliary
quantity $\norm{z_t-z_{t-1}}^{2}$ to both sides. After rearrangement we obtain
\[
\begin{aligned}
V_{t+1}
&:=\norm{z_{t+1}-z^\star}^{2}+\norm{z_{t+1}-z_t}^{2}\\
&\le \norm{z_t-z^\star}^{2}+\norm{z_t-z_{t-1}}^{2}
-\eta\bigl(1-2\eta L\bigr)\norm{F(z_t)}^{2}.
\end{aligned}
\tag{11}
\]
\noindent\textbf{[Step 1]} The inequality \eqref{11} is exactly Lemma~\ref{lem:potential}.

\noindent\textbf{Step 8.}  Since $0<\eta\le\frac{1}{4L}$, the coefficient
$1-2\eta L$ is non‑negative, hence $V_{t+1}\le V_t$. Summing \eqref{11} for
$t=0,\dots,T-1$ yields
\[
\sum_{t=0}^{T-1}\norm{F(z_t)}^{2}
\le \frac{V_0-V_T}{\eta\bigl(1-2\eta L\bigr)}
\le \frac{\norm{z_0-z^\star}^{2}}{\eta\bigl(1-2\eta L\bigr)}.
\tag{12}
\]

\noindent\textbf{Step 9.}  By convex–concavity,
\[
L(x_t,y)-L(x^\star,y_t)\le\inner{F(z_t)}{z_t-z^\star},
\qquad
\forall t.
\tag{13}
\]
Averaging \eqref{13} over $t=0,\dots,T-1$ and using Cauchy–Schwarz gives
\[
\Delta_T
\le \frac{1}{T}\sum_{t=0}^{T-1}\inner{F(z_t)}{z_t-z^\star}
\le \frac{1}{T}\sum_{t=0}^{T-1}\norm{F(z_t)}\norm{z_t-z^\star}.
\tag{14}
\]

\noindent\textbf{Step 10.}  Apply Cauchy–Schwarz again:
\[
\Delta_T\le
\frac{1}{T}\Bigl(\sum_{t=0}^{T-1}\norm{F(z_t)}^{2}\Bigr)^{\!1/2}
\Bigl(\sum_{t=0}^{T-1}\norm{z_t-z^\star}^{2}\Bigr)^{\!1/2}
\le\frac{\sqrt{T}\,\bigl(\sum_{t=0}^{T-1}\norm{F(z_t)}^{2}\bigr)^{1/2}
\max_{t}\norm{z_t-z^\star}}{T}.
\tag{15}
\]
Since $\norm{z_t-z^\star}\le \norm{z_0-z^\star}$ for all $t$ (by monotonicity of $V_t$),
\[
\Delta_T\le\frac{\norm{z_0-z^\star}}{\sqrt{T}}\,
\Bigl(\frac{\sum_{t=0}^{T-1}\norm{F(z_t)}^{2}}{T}\Bigr)^{1/2}.
\tag{16}
\]

\noindent\textbf{Step 11.}  Insert the bound \eqref{12} into \eqref{16}:
\[
\Delta_T\le
\frac{\norm{z_0-z^\star}}{\sqrt{T}}\,
\Bigl(\frac{\norm{z_0-z^\star}^{2}}{\eta T\bigl(1-2\eta L\bigr)}\Bigr)^{1/2}
=\frac{\norm{z_0-z^\star}^{2}}{\eta\bigl(1-2\eta L\bigr)T}.
\tag{17}
\]
Choosing the maximal admissible stepsize $\eta=1/(4L)$ makes
$1-2\eta L=1/2$, and we obtain the clean bound
\[
\boxed{\displaystyle
\Delta_T\le\frac{\norm{z_0-z^\star}^{2}}{\eta\,T}}.
\tag{18}
\]
This is exactly the claim \eqref{thm:ogda}. ∎

\section*{Rate Derivation (With Constants)}
From \eqref{18} we read off the explicit constants:
\[
\Delta_T\le \underbrace{\frac{\norm{z_0-z^\star}^{2}}{\eta}}_{\text{initial error term}}
\cdot\frac{1}{T}.
\]
If $\eta=1/(4L)$,
\[
\Delta_T\le 4L\,\norm{z_0-z^\star}^{2}\,\frac{1}{T}.
\]
Thus the convergence constant scales linearly with the smoothness $L$ and
quadratically with the distance of the initialization to the saddle point.

\section*{Special Cases}
\begin{enumerate}[leftmargin=*,label=\arabic*.]
\item \textbf{Strongly convex–strongly concave $L$.}
If $L$ is $\mu$‑strongly convex in $x$ and $\mu$‑strongly concave in $y$, then
$F$ is $\mu$‑strongly monotone. The same proof yields a linear rate:
\[
\Delta_T\le\bigl(1-\eta\mu\bigr)^{T}\norm{z_0-z^\star}^{2}.
\]

\item \textbf{Unconstrained Euclidean space.}
When $\mathcal X=\mathcal Y=\R^{d}$ and no projection is needed, the analysis
remains unchanged; the potential $V_t$ does not contain projection errors.

\item \textbf{Stochastic gradients.}
If $g_t$ is an unbiased estimator of $F(z_t)$ with bounded variance,
the same steps go through in expectation, adding a term
$\mathcal O(\sigma^{2}/T)$ to the right‑hand side of \eqref{thm:ogda}.
\end{enumerate}

\end{document}